\section{Digitální elektronika (BPC-DE1)}
\subsection{Logická hradla. Postuláty a zákony Booleovy algebry}
    \subsubsection{Logická hradla a základní kombinační funkce}
        Kombinační logická funkce je pravidlo, které přiřazuje každé kombinaci logických hodnot (0 a 1) vstupních proměnných z definičního oboru jedinou výstupní hodnotu. Pro daný počet vstupů $n$ existuje vždy konečný počet možných kombinačních funkcí, a to přesně $2^{2^n}$.
        
        \begin{itemize}
            \item \textbf{Funkce jedné proměnné}
            \begin{itemize}
                \item \textbf{Nulová funkce:} $y = f_0(x) = 0$.
                \item \textbf{Opakování (Buffer, BUF):} $y = f_1(x) = x$.
                \item \textbf{Inverze (NOT, negace):} $y = f_2(x) = \bar{x}$.
                \item \textbf{Jednotková funkce:} $y = f_3(x) = 1$.
            \end{itemize}
        
            \item \textbf{Funkce dvou proměnných a logická hradla}
            \begin{itemize}
                \item Ze 16 možných logických funkcí dvou proměnných se v digitální technice využívají především tři základní funkce a jejich invertované podoby:
                \item \textbf{AND (logický součin, konjunkce):} $y = a \cdot b$.
                \item \textbf{OR (logický součet, disjunkce):} $y = a + b$.
                \item \textbf{XOR (exkluzivní logický součet, nonekvivalence):} $y = a \oplus b = a \cdot \bar{b} + \bar{a} \cdot b$.
                \item Ke každé z těchto základních funkcí existuje její \textbf{negovaná (inverzní) forma}:
                \item \textbf{NAND:} $y = \overline{a \cdot b}$.
                \item \textbf{NOR:} $y = \overline{a + b}$.
                \item \textbf{XNOR (ekvivalence):} $y = \overline{a \oplus b} = a \cdot b + \bar{a} \cdot \bar{b}$.
            \end{itemize}
        
            \item \textbf{Pravdivostní stavy základních hradel pro vstupy a, b}
            \begin{itemize}
                \item \textbf{AND:} Nabývá hodnoty 1 pouze tehdy, jsou-li oba vstupy 1. V ostatních případech je 0.
                \item \textbf{OR:} Nabývá hodnoty 0 pouze tehdy, jsou-li oba vstupy 0. V ostatních případech je 1.
                \item \textbf{XOR:} Nabývá hodnoty 1, pokud jsou vstupní hodnoty rozdílné (0 a 1 nebo 1 a 0).
            \end{itemize}
        \end{itemize}
    \subsubsection{Postuláty a zákony Booleovy algebry}
        Booleova algebra představuje matematický nástroj pro manipulaci a práci s logickými (dvoustavovými) proměnnými, využívající primárně operace AND, OR a NOT. Lze dokázat, že jakákoliv logická funkce s libovolným počtem proměnných může být reprezentována (realizována) pomocí elementárních logických funkcí, jež tvoří tzv. \textbf{úplný soubor logických funkcí}.
        
        Při zjednodušování logických výrazů se aplikují následující základní postuláty a zákony:
        
        \begin{itemize}
            \item \textbf{Vliv konstant a operace identity}
            \begin{itemize}
                \item Pro konstanty 0 a 1 a identické proměnné platí:
                \item $0 + a = a$ a současně $a + 0 = a$
                \item $1 + a = 1$ a současně $a + 1 = 1$
                \item $1 \cdot a = a$ a současně $a \cdot 1 = a$
                \item $0 \cdot a = 0$ a současně $a \cdot 0 = 0$
                \item $a + a = a$
                \item $a \cdot a = a$
            \end{itemize}
            
            \item \textbf{Zákony negace}
            \begin{itemize}
                \item Pro logickou operaci negace platí pravidla komplementarity a dvojité negace:
                \item $\bar{0} = 1$ a $\bar{1} = 0$
                \item $a + \bar{a} = 1$
                \item $a \cdot \bar{a} = 0$
                \item $\bar{\bar{a}} = a$ (dvojitá negace)
            \end{itemize}

            \item \textbf{Komutativní zákon}
            \begin{itemize}
                \item Pořadí operandů u základních logických funkcí nehraje roli:
                \item $a + b = b + a$
                \item $a \cdot b = b \cdot a$
            \end{itemize}
        
            \item \textbf{Asociativní zákon}
            \begin{itemize}
                \item Operandy lze libovolně sdružovat:
                \item $a + b + c = a + (b + c)$
                \item $a \cdot b \cdot c = a \cdot (b \cdot c)$
            \end{itemize}
        
            \item \textbf{Distributivní zákon (Distributivnost)}
            \begin{itemize}
                \item Určuje pravidla roznásobování a vytýkání v logických výrazech:
                \item $a \cdot (b + c) = a \cdot b + a \cdot c$
                \item $a + b \cdot c = (a + b) \cdot (a + c)$
            \end{itemize}
        
            \item \textbf{Zákon absorpce (pohlcení)}
            \begin{itemize}
                \item Slouží ke zjednodušení logického výrazu eliminací redundantní proměnné:
                \item $a + a \cdot b = a$
                \item $a \cdot (a + b) = a$
            \end{itemize}
        
            \item \textbf{Zákon spojování}
            \begin{itemize}
                \item Využívá se při minimalizaci logických funkcí k odstranění proměnné, která se vyskytuje v přímém i negovaném tvaru:
                \item $a \cdot b + a \cdot \bar{b} = a$
                \item $(a + b) \cdot (a + \bar{b}) = a$
            \end{itemize}

            \item \textbf{Princip duality}
            \begin{itemize}
                \item Podle principu duality platí, že dva ekvivalentní logické výrazy si zachovají svou rovnost (zůstávají v platnosti) i v případě, že současně vzájemně zaměníme logické operátory AND ($\cdot$) za OR (+) a naopak, a k tomu zaměníme všechny logické konstanty 0 za 1 a naopak.
            \end{itemize}
        
            \item \textbf{De Morganovy zákony}
            \begin{itemize}
                \item Základní věty definující chování negace u logických operací. Hodnota logického výrazu se nezmění, pokud operátory logického součtu nahradíme operátory logického součinu (a naopak), a současně invertujeme všechny proměnné i celý výsledek.
                \item Vyjadřují se dvěma základními teorémy:
                \item \textbf{Teorém 1:} $\overline{a \cdot b} = \bar{a} + \bar{b}$
                \item \textbf{Teorém 2:} $\overline{a + b} = \bar{a} \cdot \bar{b}$
            \end{itemize}
        \end{itemize}
\subsection{Klopné obvody. Typy, vlastnosti, zapojení a realizace.}
        Klopné obvody tvoří základní paměťové prvky v sekvenčních logických obvodech. Na rozdíl od kombinačních obvodů závisí jejich výstupní hodnoty nejen na aktuálním stavu vstupních signálů, ale také na předchozích hodnotách (obsahují paměťový segment).
        
        Podle způsobu řízení a reakce na vstupní signály se klopné obvody dělí na dvě hlavní kategorie:
        \begin{itemize}
            \item \textbf{Latches (Klopné obvody se statickým řízením):} Zařízení citlivá na úroveň řídicího signálu (\textit{level-sensitive}). Obvod je po dobu aktivní úrovně transparentní.
            \item \textbf{Flip-flops (Klopné obvody s dynamickým řízením):} Zařízení citlivá výhradně na hranu řídicího (hodinového) signálu (\textit{edge-triggered}). Výstup reaguje pouze v okamžiku náběžné nebo sestupné hrany.
        \end{itemize}
    \subsubsection{Latches: Klopné obvody citlivé na úroveň}
        \begin{itemize}
            \item \textbf{S-R Latch}
            \begin{itemize}
                \item \textbf{Zapojení a realizace:} Typicky se realizuje asymetrickým křížovým spojením dvou dvouvstupových hradel NOR nebo NAND.
                \item \textbf{Vlastnosti (realizace s hradly NOR):} Má dva vstupy s (set – nastavení) a r (reset – nulování). Logická "1" na vstupu s nastaví hlavní výstup do stavu 1. Logická "1" na vstupu r výstup vynuluje. Stav, kdy s=0 a r=0, je paměťový (obvod drží předchozí hodnotu). Kombinace s=1, r=1 je \textbf{nepovolená} – výstupy ztrácí komplementaritu (Q i /Q jsou obojí 0). Při návratu z tohoto zakázaného stavu zpět do stavu paměťového (00) se obvod na desítky nanosekund dostává do tzv. metastabilního stavu s nedefinovaným výstupem.
                \item \textbf{Vlastnosti (realizace s hradly NAND):} Vstupy /S a /R jsou aktivní v nízké úrovni. Paměťový stav je pro "11" a zakázaný (nepovolený) naopak pro hodnoty "00".
            \end{itemize}
        
            \item \textbf{Gated S-R Latch}
            \begin{itemize}
                \item \textbf{Zapojení a realizace:} Vznikne přidáním povolovacího vstupu \textit{en} (enable) před základní S-R obvod. Ve struktuře složitějších integrovaných obvodů se využívá velmi často jako dílčí stavební blok.
                \item \textbf{Vlastnosti:} Obvod reaguje na vstupy s a r, pouze pokud je povolovací vstup v aktivní úrovni (např. \textit{en}=1). Pokud je \textit{en}=0, vstupy pro nastavení a reset jsou deaktivovány, obvod si pamatuje svou hodnotu a jakékoliv změny datových vstupů jsou na výstupu ignorovány.
            \end{itemize}
        
            \item \textbf{D Latch}
            \begin{itemize}
                \item \textbf{Zapojení a realizace:} Vychází z obvodu Gated S-R, ke kterému je přidán invertor. Tím je zaručeno, že hodnoty na větvích s a r jsou neustále logickým komplementem. Původní vstup "s" je přejmenován na D (data).
                \item \textbf{Vlastnosti:} Implementace invertoru eliminuje možnost vzniku zakázaného vstupního stavu. Když \textit{en}=1, je obvod transparentní (výstup Q přímo kopíruje vstup D). Po přechodu \textit{en} na 0 se na sestupné hraně aktuální hodnota zablokuje a zapamatuje do příchodu dalšího povolovacího pulzu.
            \end{itemize}
        \end{itemize}
        \newpage
        \begin{figure}[h]
            \centering
            \includegraphics[width=1\linewidth]{SRlatch.png}
            \caption{S-R Latch}
        \end{figure}
        \begin{figure}[h]
            \centering
            \includegraphics[width=1\linewidth]{GatedSRlatch.png}
            \caption{Gated S-R Latch}
        \end{figure}
        \begin{figure}[h]
            \centering
            \includegraphics[width=1\linewidth]{Dlatch.png}
            \caption{D Latch}
        \end{figure}
    
    \subsubsection{Flip-Flops: Klopné obvody řízené hranou}

        Pro tyto obvody je typická přítomnost hodinového signálu (\textit{clk}). Citlivost na hranu je ve schematické značce u pinu \textit{clk} znázorněna pomocí malého trojúhelníku.
        
        \begin{itemize}
            \item \textbf{D-type Flip-Flop}
            \begin{itemize}
                \item \textbf{Zapojení a realizace:} Vytváří se v kaskádové topologii Master-Slave propojením dvou klopných obvodů: D latche a Gated S-R latche, které jsou aktivovány v opačných fázích hodinového signálu.
                \item \textbf{Vlastnosti:} Během první fáze hodin propisuje první obvod (Master) data na svůj výstup, zatímco druhý (Slave) je deaktivovaný. Teprve při aktivní hraně hodin (např. sestupné) se aktivuje stupeň Slave a výstupní hodnota prvního se překlopí na konečný výstup Q.
                \item \textbf{Charakteristická rovnice:} Výstupní stav po aktivní hraně je definován rovnicí $q_{n+1} = d$.
            \end{itemize}
        
            \item \textbf{JK-type Flip-Flop}
            \begin{itemize}
                \item \textbf{Zapojení a realizace:} Modifikovaný hranově řízený D-type Master-Slave obvod, kde jsou výstupy klopného obvodu přivedeny zpět na vstup a logicky zkombinovány s datovými vstupy j a k.
                \item \textbf{Vlastnosti:} Obvod řeší problém nedefinovaného stavu klasického RS obvodu. Při kombinaci vstupů j=1 a k=1 se namísto chybového stavu výstup Q invertuje (tzv. toggle) vůči předchozímu stavu. Zbylé logické stavy jsou totožné s RS (00 = paměť, 01 = reset do nuly, 10 = nastavení jedničky).
                \item \textbf{Charakteristická rovnice:} $q_{n+1} = j \cdot \bar{q}_n + \bar{k} \cdot q_n$.
            \end{itemize}
        
            \item \textbf{T-type Flip-Flop}
            \begin{itemize}
                \item \textbf{Zapojení a realizace:} Je strukturálně identický s obvodem JK, se zásadním rozdílem – datové vstupy j a k jsou zkratovány do jediného datového vstupu t.
                \item \textbf{Vlastnosti:} Působí jako obvod pro invertování výstupu. Pokud je na datovém vstupu t nula, obvod svůj stav udržuje. Je-li přítomna jedna, na aktivní hraně hodin se aktuální uložený stav neguje.
                \item \textbf{Charakteristická rovnice:} $q_{n+1} = t \cdot \bar{q}_n + \bar{t} \cdot q_n = t \oplus q_n$.
            \end{itemize}
        \end{itemize}
    \subsubsection{Časové parametry}
        Pro bezchybnou realizaci záznamu dat u synchronních překlápěcích obvodů je nutné bezpodmínečně respektovat fyzikální zpoždění součástek v podobě specifických časových oken. Jejich narušení u datových signálů může způsobit zapsání chybné hodnoty či metastabilitu.
        
        \begin{itemize}
            \item \textbf{Setup time ($t_{su}$):} Určuje minimální přesný časový okamžik před příchodem aktivní hrany hodin, po který již musí být hodnota na datovém vstupu neměnná (stabilní).
            \item \textbf{Hold time ($t_h$):} Představuje minimální nutný čas bezprostředně po aktivní hraně hodin, během kterého se zapisovaná data stále nesmí změnit.
            \item \textbf{Clock-to-Q delay ($t_{cq}$):} Doba nutná pro propagaci zpracovávaného signálu, reprezentující prodlevu od reálného příchodu hrany hodinových pulzů do stabilizování nové hodnoty přímo na fyzickém výstupu Q.
        \end{itemize}
\subsection{Čítače. Typy, vlastnosti, zapojení a realizace. }
        Čítač je sekvenční logický obvod, který má schopnost počítat vstupní hodinové impulzy. Nejběžnější typ čítače má jeden hodinový vstup a více výstupů, které zpravidla reprezentují binární nebo binárně kódovaná desetinná (BCD) čísla.
        
        \begin{itemize}
            \item \textbf{Základní vlastnosti čítačů:}
            \begin{itemize}
                \item \textbf{Maximální počet stavů:} Obecný čítač s $n$ klopnými obvody může mít maximálně $2^n$ stavů.
                \item \textbf{Modul (MOD):} Označuje počet výstupních stavů, kterými čítač projde ve výstupní sekvenci, než se vrátí ke své počáteční hodnotě. Například čítač sestavený ze 4 klopných obvodů má maximum $2^4 = 16$ stavů (0--15) a označuje se jako "MOD-16" čítač.
            \end{itemize}
        \end{itemize}
    \subsubsection{Typy a klasifikace čítačů}
        \begin{itemize}
            \item \textbf{Podle vnitřního propojení (synchronizace):} Asynchronní a synchronní.
            \item \textbf{Podle kódu:} Binární, dekadické (\textit{decade}), čítače v Grayově kódu, LFSR (\textit{Linear Feedback Shift Register}) a další.
            \item \textbf{Podle směru čítání:} Čítající vpřed (\textit{UP counter}), čítající vzad (\textit{DOWN counter}) a vratné (\textit{UP/DOWN counter}).
            \item \textbf{Podle typu použitých klopných obvodů:} Typicky D, T nebo JK.
        \end{itemize}
    \subsubsection{Asynchronní čítače (Ripple counters)}
        U asynchronních čítačů jsou jednotlivé klopné obvody (KO) řízeny různými hodinovými signály, výstupy tak mění svůj stav jeden po druhém.
        \begin{itemize}
            \item \textbf{Zapojení a realizace}
            \begin{itemize}
                \item Externím hodinovým signálem \textit{clk} je taktován pouze první klopný obvod. Hodinový vstup každého dalšího stupně je tvořen výstupem z předchozího klopného obvodu.
                \item \textbf{Čítač vpřed (UP):} Realizován s D klopnými obvody tak, že vstup D každého KO je připojen na jeho vlastní invertovaný výstup ($\bar{Q}$). Tento invertovaný výstup slouží zároveň jako hodinový signál pro následující klopný obvod.
                \item \textbf{Čítač vzad (DOWN):} Lze realizovat totožným zapojením, jedinou změnou je použití hlavního neinvertovaného výstupu ($Q$) pro hodinový vstup následujícího klopného obvodu.
                \item Nevyžadují žádná dodatečná logická hradla, zapojení je velmi jednoduché.
            \end{itemize}
        
            \item \textbf{Vlastnosti a omezení}
            \begin{itemize}
                \item \textbf{Ripple effect:} Při přechodu (např. ze stavu 111 na 000) se výstupy mění sekvenčně. Zpoždění klopných obvodů způsobuje, že se v krátkém časovém úseku na výstupu generují nepravé hodnoty.
                \item \textbf{Maximální frekvence:} Doba ustálení je přímo úměrná počtu bitů čítače $n$. Maximální frekvence hodin se vypočítá jako $$f_{clkmax} = \frac{1}{n \cdot t_{FF} + t_{setup}}$$ kde $t_{FF}$ je zpoždění klopného obvodu a $t_{setup}$ je doba předstihu.
            \end{itemize}
        \end{itemize}
    \subsubsection{Synchronní čítače}

        V synchronním čítači je externí hodinový signál (\textit{clk}) připojen k hodinovým vstupům všech klopných obvodů současně. Všechny KO jsou tedy taktovány paralelně ve stejném čase.
        
        \begin{itemize}
            \item \textbf{Vlastnosti a zapojení}
            \begin{itemize}
                \item Klopné obvody tvoří registr, a protože mění stav současně, nevzniká zde žádný "ripple effect".
                \item Kromě registru obsahuje synchronní systém kombinační obvody ve zpětnovazební smyčce, které definují následující stav.
            \end{itemize}
        
            \item \textbf{Binární synchronní čítač}
            \begin{itemize}
                \item Z pravdivostní tabulky vyplývá, že nejnižší bit (LSB) se invertuje s každou aktivní hranou hodin. Každý další vyšší bit se invertuje pouze tehdy, jsou-li všechny aktuální nižší bity rovny 1.
                \item \textbf{Realizace s T klopnými obvody:} Pro inverzi stavu využívá KO typu T. K ověření, že jsou všechny nižší bity rovny 1, se využívají hradla AND.
                \item \textbf{Varianty přenosu:}
                \begin{enumerate}
                    \item \textbf{Paralelní přenos (Parallel carry):} Využívá hradla AND s rostoucím počtem vstupů připojená přímo na výstupy předchozích KO. Maximální frekvence: $$f_{clkmax} = \frac{1}{t_{FF} + t_{AND} + t_{setup}}$$
                    \item \textbf{Sériový přenos (Ripple carry):} Místo rozšiřování počtu vstupů hradel AND se využívá řetězení 2vstupových hradel AND. Tímto způsobem jsou čítače implementovány v programovatelných logických obvodech (PLD). Maximální frekvence: $$f_{clkmax} = \frac{1}{t_{FF} + (n-1) \cdot t_{AND} + t_{setup}}$$
                \end{enumerate}
                \item \textbf{Realizace s D klopnými obvody:} Protože pro převod T klopného obvodu na D platí $d = t \oplus q_n$, objevuje se v kombinační části obvodu kromě hradel AND také hradlo XOR.
            \end{itemize}
        
            \item \textbf{LFSR čítač (Linear Feedback Shift Register)}
            \begin{itemize}
                \item Implementován jako posuvný registr tvořený sérií klopných obvodů. Zpětná vazba pro začátek řetězce je tvořena výstupem hradla XOR nebo XNOR, do kterého vstupují vybrané výstupy registru (taps).
                \item Generuje pseudonáhodnou binární sekvenci. Maximální počet stavů $n$-bitového LFSR není $2^n$, ale $2^n - 1$. Důvodem je, že např. při použití hradla XOR by sekvence samých nul trvale produkovala další nuly. Dosažení maximální délky je podmíněno správným výběrem zpětnovazebních odboček.
                \item \textbf{Výhody:} Jednoduchá vnitřní struktura a vysoká rychlost.
            \end{itemize}
        \end{itemize}